%% ma: L12-B13-0029
%% lop: 12
%% chuong: 4
%% bai: 13
%% dang: 12.13.12
%% loai: TLN
%% muc: VD
%% dap_an: "192"
%% nguon: "(NBV)-ĐỀ SỐ 13-MỖI NGÀY 1 ĐỀ THI 2025.pdf (D13, MỖI NGÀY 1 ĐỀ - Đề số 13), Phần III Câu 6"
%% kien_thuc: "Bài 13 (thể tích khối tròn xoay), Bài 12 (tích phân), Bài 11 (nguyên hàm)"
%% trang_thai: chinh-thuc
%% ghi_chu: "Đáp án gốc 192 đúng (sympy: parabol y = 50 - x^2/250 quay quanh trục hoành, thể tích mỗi thùng 656000 pi/3 cm^3 khoảng 686,96 lít; 7 thùng 4808,7 lít, tiền khoảng 192,35 triệu đồng). Cùng ý tưởng (đường sinh parabol, khối tròn xoay) với dạng 12.13.12 nhưng quay quanh trục hoành. Đã bỏ ảnh thùng rượu minh họa (không chứa dữ kiện toán)."
\begin{ex}
Để chuẩn bị cho bữa tiệc cuối năm, công ty của ông Nam có đặt mua từ cơ sở sản xuất $7$ thùng rượu với kích thước như nhau, thùng có dạng khối tròn xoay với đường sinh dạng parabol, mỗi thùng rượu có bán kính hai mặt là $40$ cm và ở giữa là $50$ cm. Chiều dài mỗi thùng rượu là $100$ cm. Biết rằng thùng rượu chứa đầy rượu và giá mỗi lít rượu là $40\,000$ đồng. Tính số tiền (đơn vị: triệu đồng) mà công ty của ông Nam phải trả cho cơ sở sản xuất rượu (kết quả làm tròn đến hàng đơn vị).
\shortans{192}
\loigiai{Đặt trục $Ox$ dọc theo chiều dài thùng, gốc $O$ ở tâm thùng. Đường sinh là parabol $y=ax^2+50$ đi qua $(50;40)$ nên $2500a+50=40$, $a=-\dfrac1{250}$.

Thể tích một thùng: $V_1=\pi\displaystyle\int_{-50}^{50}\left(50-\dfrac{x^2}{250}\right)^2\mathrm dx=\dfrac{656000\pi}3\ (\text{cm}^3)\approx686{,}96$ lít.

Số tiền cho $7$ thùng: $7\cdot686{,}96\cdot40\,000\approx192{,}3\cdot10^6$ đồng, tức khoảng $192$ triệu đồng.}
\end{ex}
